% 1-4-contributions.tex
%  Budget: 1pp.
%  Rules: every numeral in \lr{}; cite only from references.bib;
%  one unit at a time -- see thesis/WRITING_HANDOFF.md and APPROVALS.md.

\section{نوآوری‌ها و دستاوردهای پژوهش}

دستاوردهای این پژوهش را می‌توان در پنج محور خلاصه کرد. تأکید همه آن‌ها بر
\emph{اعتبار مقایسه} است، نه بر افزودن یک معماری تازه به ادبیات موجود.

\textbf{یکم، مقایسه‌ای که به ادبیات منتشرشده گره خورده است.} ارزیابی بر
مجموعه‌داده عمومی \lr{EPEX-DE} و مطابق پروتکل پیشنهادی لاگو و همکاران
\cite{lago_forecasting_2021} انجام شده است. این انتخاب امکان مقایسه مستقیم
نتایج این پژوهش با اعداد منتشرشده همان مرجع را فراهم می‌کند \lr{---} امکانی که در
بیشتر مطالعات این حوزه، به دلیل استفاده از داده‌های منتشرنشده، وجود ندارد.

\textbf{دوم، عملیاتی‌کردن همزمان دو افق.} هدف ساعتی (بردار \lr{24}‌بعدی) و
هدف روزانه (میانگین بار پایه) هر دو پیاده‌سازی شده‌اند، و هدف روزانه از دو
مسیر مستقل \lr{---} مدل‌سازی مستقیم و تجمیع پیش‌بینی‌های ساعتی \lr{---} محاسبه شده
است. مقایسه این دو مسیر خود پاسخ پرسش چهارم پژوهش است.

\textbf{سوم، آزمون معناداری آماری به‌جای مقایسه میانگین‌ها.} تفاوت عملکرد
مدل‌ها با آزمون دیبولد-ماریانو و با تصحیح \lr{HAC} سنجیده شده و برای خانواده
آزمون‌های چندگانه نیز تصحیح شده است. نتیجه عملی این تصحیح آن است که برخی
تفاوت‌هایی که در مقایسه خام معنادار به نظر می‌رسیدند، پس از تصحیح معنادار
نیستند؛ این موارد صریحاً گزارش شده‌اند.

\textbf{چهارم، وزن‌دهی مشروط به رژیم بازار.} ترکیب مدل‌ها نه با وزن ثابت،
بلکه با دو مجموعه وزن مجزا برای شرایط آرام و پرتنش انجام شده است. آستانه
تفکیک این دو رژیم صرفاً از داده اعتبارسنجی استخراج شده و هرگز بر پایه رفتار
داده آزمون تنظیم نشده است.

\textbf{پنجم، و از نظر روش‌شناختی مهم‌ترین دستاورد: ارزیابی برون‌توزیع و
حسابرسی آن.} مدل‌های تثبیت‌شده بر دوره بنچمارک، بدون هیچ بازآموزی، بر داده
زنده سال \lr{2026} آزموده شده‌اند. نتیجه منفی است \lr{---} که خود یافته‌ای
درخور گزارش است \lr{---} اما این پژوهش به گزارش نتیجه منفی بسنده نکرده و سازوکار
آن را نیز سنجیده است. چنان‌که در فصل پنجم نشان داده می‌شود، بخش قابل
اندازه‌گیری از این افت، ناشی از عدم انطباق دامنه ویژگی‌های برون‌زا میان دوره
آموزش و دوره سرویس‌دهی است و نه صرفاً ناشی از تغییر بازار. تفکیک این دو
سازوکار، ادعای پژوهش را محدودتر و در عین حال دفاع‌پذیرتر می‌کند.

افزون بر این پنج محور، خروجی نرم‌افزاری پژوهش نیز به‌صورت یک ابزار
پیش‌بینی عملیاتی تدوین شده است که در بخش \lr{5-3} معرفی می‌شود.
